% Publication-ready vector figure: instructional mechanism
\begin{figure}[htbp]
\centering
\begin{tikzpicture}[
  font=\small,
  node distance=4mm,
  stage/.style={draw, rounded corners=2pt, align=center, minimum height=8mm, text width=110mm, inner sep=4pt},
  arrow/.style={-{Latex[length=2.2mm]}, line width=0.7pt}
]
\node[stage] (predict) {\textbf{Predict}\\Anticipate how a parameter change should affect the signal or result};
\node[stage, below=of predict] (manipulate) {\textbf{Manipulate}\\Change one DSP parameter or select a comparison case};
\node[stage, below=of manipulate] (observe) {\textbf{Observe}\\Inspect immediate waveform, spectrum, or numerical feedback};
\node[stage, below=of observe] (compare) {\textbf{Compare}\\Contrast parameter settings or guided cases};
\node[stage, below=of compare] (explain) {\textbf{Explain and reflect}\\Articulate the DSP relationship using the observed evidence};

\draw[arrow] (predict) -- (manipulate);
\draw[arrow] (manipulate) -- (observe);
\draw[arrow] (observe) -- (compare);
\draw[arrow] (compare) -- (explain);
\draw[arrow] (explain.east) -- ++(8mm,0) |- (predict.east);
\end{tikzpicture}
\caption{Instructional cycle embedded in the DSP-Lab activities. The design emphasizes guided conceptual comparison rather than unrestricted simulation.}
\label{fig:learning-cycle}
\end{figure}
